%%% =========================================================================
%%  Factsheet — Template (classe DPQISTfactsheet)
%%  Direção de Planeamento e Qualidade — IST
%%  Compilar com XeLaTeX ou LuaLaTeX
%%% =========================================================================

\documentclass[a4paper, 11pt]{DPQISTfactsheet}

%----------------------------------------------------------------------------------------
%	DEFINIÇÕES DO FACTSHEET
%----------------------------------------------------------------------------------------

\setfactsheettitle{Factsheet — Título do Indicador}
\setfactsheetsubtitle{Subtítulo informativo curto}
\setfactsheetcode{FS-DPQ-001}           % código interno do documento
\setfactsheetdate{\today}
\setdocversion{v1.00}

%----------------------------------------------------------------------------------------
\begin{document}

%----------------------------------------------------------------------------------------
%	TÍTULO DO FACTSHEET (bloco de abertura)
%----------------------------------------------------------------------------------------

\begin{center}
    {\fontsize{26pt}{30pt}\selectfont\bfseries\color{dpqheader}\factsheettitle\par}
    \vspace{0.3cm}
    {\Large\itshape\factsheetsubtitle\par}
    \vspace{0.2cm}
    {\small\color{istdarkgray}\docversion\ · \factsheetdate\ · \factsheetcode\par}
\end{center}

\vspace{0.4cm}

%----------------------------------------------------------------------------------------
%	INDICADORES-CHAVE (key metrics) — opcional
%----------------------------------------------------------------------------------------

\begin{tabularx}{\textwidth}{XXX}
    \begin{keymetricbox}
        \keylabel{Indicador 1}
        \keyvalue{\SI{42,3}{\percent}}
        \keynote{Variação de +2,1 p.p. face a 2025}
    \end{keymetricbox}
    &
    \begin{keymetricbox}
        \keylabel{Indicador 2}
        \keyvalue{\num{12345}}
        \keynote{Fonte: base de dados institucional}
    \end{keymetricbox}
    &
    \begin{keymetricbox}
        \keylabel{Indicador 3}
        \keyvalue{\SI{97,8}{\percent}}
        \keynote{Meta anual: \SI{95}{\percent}}
    \end{keymetricbox}
\end{tabularx}

%----------------------------------------------------------------------------------------
%	SECÇÃO
%----------------------------------------------------------------------------------------

\section{Enquadramento}

Texto introdutório do factsheet. Sem capa nem registo de alterações: o documento
começa directamente pelo título e pelos resultados.

%----------------------------------------------------------------------------------------
%	TABELA NUMÉRICA PEQUENA/MÉDIA (portrait)
%----------------------------------------------------------------------------------------

\subsection{Resultados por categoria}

As colunas \verb|S| alinham os números à direita pela vírgula decimal e
agrupam os milhares automaticamente.

\begin{table}[htbp]
    \centering
    \caption{Exemplo de tabela numérica em portrait}
    \label{tab:portrait}
    \rowcolors{2}{dpqrowlight}{white}
    \begin{tabular}{L{4cm} S[table-format=5.0] S[table-format=2.1] S[table-format=2.1]}
        \toprule
        {\theadS{Categoria}} & {\theadS{N}} & {\theadS{Média}} & {\theadS{Desvio padrão}} \\
        \midrule
        Categoria A & 12345 & 87,3 & 4,2 \\
        Categoria B & 9870  & 82,1 & 5,7 \\
        Categoria C & 6540  & 78,9 & 6,1 \\
        \bottomrule
    \end{tabular}
\end{table}

%----------------------------------------------------------------------------------------
%	TABELA NUMÉRICA GRANDE EM LANDSCAPE
%----------------------------------------------------------------------------------------

\newpage

\begin{landscape}

\section{Dados detalhados (landscape)}

Em landscape a largura útil aumenta, permitindo tabelas de 8–12 colunas.
O cabeçalho e rodapé acompanham a rotação da página.

\begin{table}[htbp]
    \centering
    \caption{Exemplo de tabela numérica grande em landscape}
    \label{tab:landscape}
    \rowcolors{2}{dpqrowlight}{white}
    \resizebox{\linewidth}{!}{%
    \begin{tabular}{L{2.5cm} *{8}{S[table-format=5.1]}}
        \toprule
        {\theadS{Unidade orgânica}} & {\theadS{2023}} & {\theadS{2024}} & {\theadS{2025}} & {\theadS{2026}} &
        {\theadS{$\Delta$ 24/23}} & {\theadS{$\Delta$ 25/24}} & {\theadS{$\Delta$ 26/25}} & {\theadS{Meta 2027}} \\
        \midrule
        Unidade A & 103,2 & 108,5 & 111,7 & 114,2 & 5,3 & 3,2 & 2,5 & 120,0 \\
        Unidade B & 96,8  & 99,1  & 101,3 & 102,0 & 2,3 & 2,2 & 0,7 & 105,0 \\
        Unidade C & 88,4  & 91,0  & 95,2  & 97,7  & 2,6 & 4,2 & 2,5 & 100,0 \\
        \addlinespace
        \textbf{Total} & \textbf{288,4} & \textbf{298,6} & \textbf{308,2} & \textbf{313,9} & \textbf{10,2} & \textbf{9,6} & \textbf{5,7} & \textbf{325,0} \\
        \bottomrule
    \end{tabular}%
    }
\end{table}

Notas de utilização: use \verb|\resizebox{\linewidth}{!}| para compressão
automática, \verb|\addlinespace| + \verb|\textbf| para linhas de total, e
\verb|\theadS{...}| em todas as células da primeira linha (fundo azul, texto branco).

\end{landscape}

%----------------------------------------------------------------------------------------
%	CAIXAS INSTITUCIONAIS (opcional)
%----------------------------------------------------------------------------------------

\newpage
\begin{istobj}
Sintaxe idêntica ao modelo de relatórios, para consistência visual entre documentos.
\end{istobj}

\begin{istprin}
\begin{enumerate}[label=\alph*), leftmargin=*]
    \item Principais conclusões em formato compacto.
    \item Fontes citadas em nota de rodapé.\footnote{Fonte: DPQ, 2026.}
\end{enumerate}
\end{istprin}

\end{document}
